\section{Retain the 129th bit of a modular sum}
Now build the smaller operation used by modular multiplication. Its inputs satisfy
\(0\leq A,B<N\), so \(S=A+B<2N\). The result is either \(S\), when \(S<N\), or
\(S-N\). Only one subtraction can be necessary, but \(S\) may require 129 bits.
Store its low 128 bits and keep the final carry separately.

Within each limb, add the two zero-extended words and the previous carry in a 64-bit
register. This temporary needs at most 33 bits. Store its low word, then shift right
by 32 to obtain the next carry. Load both operand words before storing the result
word so that an output exactly equal to either operand remains valid.

\Needspace{11\baselineskip}
\MARSCodeLabel{Produce one low result limb and its one-bit carry.}
\begin{minted}[samepage=true]{mars-idea}
daddu $t3, $t3, $t4
    daddu $t3, $t3, $t1
    daddu $t2, $a3, $t0
    sw $t3, 0($t2)
    dsrl $t1, $t3, 32
\end{minted}

If the final carry is zero, compare the stored result with \(N\) from the highest
limb down. Subtract when it is greater than or equal to \(N\). Equality matters:
\(N\bmod N\) is zero. If the carry is one, subtraction is required even when the
stored low 128 bits compare below \(N\).

For a concrete carry case, let \(N=2^{128}-159\) and \(A=B=N-1\). The true sum is
\(2^{129}-320\). Its low 128 bits are \(2^{128}-320\), which are below \(N\).
Returning those bits would be wrong. The carried sum minus \(N\) is
\(2^{128}-161=N-2\), the correct residue. Subtracting \(N\) from the stored low
bits with wraparound produces those same final 128 bits; the retained carry explains
why this subtraction must happen.

Subtraction needs a borrow between limbs. For zero-extended limbs \(z_i,n_i\) and
incoming borrow \(b\), compute \(D=z_i-n_i-b\). Its signed mathematical range is
\([-2^{32},2^{32}-1]\), which fits in a signed 64-bit temporary. A negative result
has bit 63 set, so a logical shift by 63 extracts the outgoing borrow as zero or one.
This use of the sign bit relies on that range; it is not a general unsigned 64-bit
borrow test.

\Needspace{10\baselineskip}
\MARSCodeLabel{Keep the borrow before advancing to the next limb.}
\begin{minted}[samepage=true]{mars-idea}
dsubu $t3, $t3, $t4
    dsubu $t3, $t3, $t1
    dsrl $t1, $t3, 63
    sw $t3, 0($t2)
\end{minted}

Reject a zero modulus before reading operands or storing any low limb. This guard is
part of the function boundary, not a test to defer until subtraction. For modulus one,
the only permitted residue inputs to this addition are both zero.
